\documentclass[a4paper,12pt]{article}
\usepackage[T2A]{fontenc}
\usepackage[utf8]{inputenc}
\usepackage[russian]{babel}
\usepackage{amsmath, amssymb}
\usepackage{geometry}
\geometry{top=2cm, bottom=2cm, left=2.5cm, right=1.5cm}

\begin{document}

\title{Ответы на контрольные вопросы \\
Лабораторная работа №6 \\
«Численное решение обыкновенных дифференциальных уравнений»}
\author{}
\date{}
\maketitle

\begin{enumerate}

\item \textbf{Сформулируйте задачу Коши для дифференциального уравнения 1 порядка.}

Задача Коши для обыкновенного дифференциального уравнения (ОДУ) первого порядка заключается в нахождении функции \( y(x) \), удовлетворяющей дифференциальному уравнению
\[
y' = f(x, y)
\]
и начальному условию
\[
y(x_0) = y_0,
\]
где \( f(x, y) \) — заданная функция, \( x_0 \) и \( y_0 \) — заданные числа (начальные данные). Геометрически это означает поиск интегральной кривой, проходящей через заданную точку \( (x_0, y_0) \) на плоскости \( (x, y) \).

\item \textbf{Что является решением для дифференциального уравнения 1 порядка?}

Решением ОДУ первого порядка называется функция \( y = \varphi(x) \), определённая и дифференцируемая на некотором интервале \( (a, b) \), которая при подстановке в уравнение обращает его в тождество:
\[
\varphi'(x) \equiv f(x, \varphi(x)) \quad \forall x \in (a, b).
\]
В общем случае решение содержит одну произвольную постоянную \( C \), которая определяется из начального условия.

\item \textbf{В чем заключается суть метода конечных разностей?}

Метод конечных разностей (разностный метод) — это численный подход к решению дифференциальных уравнений, при котором непрерывная область изменения независимой переменной заменяется дискретной сеткой узлов \( x_0, x_1, \dots, x_n \). Производные аппроксимируются разностными отношениями (конечными разностями), а дифференциальное уравнение заменяется системой алгебраических уравнений относительно значений искомой функции в узлах сетки.

\item \textbf{Что такое разностная аппроксимация?}

Разностная аппроксимация — это замена производной функции \( y'(x) \) приближённым выражением через значения функции в соседних узлах сетки. Например:
\begin{itemize}
    \item Правая разность: \( y'(x_i) \approx \frac{y_{i+1} - y_i}{h} \)
    \item Левая разность: \( y'(x_i) \approx \frac{y_i - y_{i-1}}{h} \)
    \item Центральная разность: \( y'(x_i) \approx \frac{y_{i+1} - y_{i-1}}{2h} \)
\end{itemize}
Здесь \( h = x_{i+1} - x_i \) — шаг сетки. Погрешность аппроксимации для центральной разности имеет порядок \( O(h^2) \).

\item \textbf{Геометрический смысл задачи Коши?}

Геометрически задача Коши для уравнения \( y' = f(x, y) \) означает: в каждой точке \( (x, y) \) области определения задано направление касательной к интегральной кривой, определяемое значением \( f(x, y) \). Требуется найти кривую, проходящую через начальную точку \( (x_0, y_0) \), у которой касательная в каждой точке совпадает с этим заданным направлением.

\item \textbf{Что такое интегральная кривая?}

Интегральная кривая — это график решения \( y = \varphi(x) \) дифференциального уравнения в плоскости \( (x, y) \). Через каждую точку области, где выполнены условия теоремы существования и единственности, проходит ровно одна интегральная кривая.

\item \textbf{Какое из условий теоремы существования и единственности решения задачи Коши для ОДУ является условием существования и какое условием единственности?}

Для уравнения \( y' = f(x, y) \) с начальным условием \( y(x_0) = y_0 \):
\begin{itemize}
    \item \textbf{Существование} обеспечивает \textbf{непрерывность} функции \( f(x, y) \) в некоторой окрестности точки \( (x_0, y_0) \).
    \item \textbf{Единственность} обеспечивает \textbf{выполнение условия Липшица} по переменной \( y \) в этой окрестности: существует константа \( L > 0 \) такая, что
    \[
    |f(x, y_1) - f(x, y_2)| \le L |y_1 - y_2|.
    \]
    Достаточным условием для Липшицевости является существование и ограниченность частной производной \( \frac{\partial f}{\partial y} \).
\end{itemize}

\item \textbf{Что должно быть задано для решения ОДУ приближенными методами?}

Для численного решения задачи Коши необходимо задать:
\begin{enumerate}
    \item Функцию \( f(x, y) \) (правую часть ОДУ).
    \item Начальные условия \( x_0, y_0 \).
    \item Интервал интегрирования \( [x_0, x_n] \).
    \item Шаг сетки \( h \) (или количество узлов).
    \item Требуемую точность \( \varepsilon \) (для методов с автоматическим выбором шага).
    \item Выбор численного метода.
\end{enumerate}

\item \textbf{Какой порядок точности имеет метод Эйлера? Рунге-Кутта?}

\begin{itemize}
    \item \textbf{Метод Эйлера} (явный) имеет \textbf{первый порядок точности}: локальная погрешность \( O(h^2) \), глобальная \( O(h) \).
    \item \textbf{Усовершенствованный метод Эйлера} (метод Эйлера-Коши, метод предиктор-корректор 2-го порядка) имеет \textbf{второй порядок точности}: глобальная погрешность \( O(h^2) \).
    \item \textbf{Классический метод Рунге-Кутта 4-го порядка} имеет \textbf{четвёртый порядок точности}: глобальная погрешность \( O(h^4) \).
\end{itemize}

\item \textbf{Перечислите основные одношаговые методы для численного решения ОДУ.}

Одношаговые методы используют информацию только в текущем узле \( (x_i, y_i) \) для вычисления \( y_{i+1} \):
\begin{itemize}
    \item Метод Эйлера (1-й порядок)
    \item Неявный метод Эйлера (1-й порядок)
    \item Усовершенствованный метод Эйлера (метод Эйлера-Коши, 2-й порядок)
    \item Модифицированный метод Эйлера (метод средней точки, 2-й порядок)
    \item Методы Рунге-Кутты 2-го, 3-го и 4-го порядков
    \item Метод Рунге-Кутты-Фельберга (с автоматическим контролем шага)
\end{itemize}

\item \textbf{Перечислите основные многошаговые методы для численного решения ОДУ.}

Многошаговые методы используют значения решения в нескольких предыдущих узлах:
\begin{itemize}
    \item \textbf{Метод Адамса} (семейство: Адамса-Башфорта — явные, Адамса-Мултона — неявные)
    \item \textbf{Метод Милна} (4-шаговый метод, часто используется как предиктор-корректор)
    \item Метод Хэмминга
    \item Неявные методы дифференцирования назад (BDF)
\end{itemize}

\item \textbf{В чем заключается суть методов прогноза и коррекции?}

Суть методов предиктор-корректор:
\begin{enumerate}
    \item \textbf{Прогноз (предиктор)}: по явной формуле (например, метод Адамса-Башфорта) вычисляется приближённое значение \( y_{i+1}^{(0)} \) в новом узле.
    \item \textbf{Коррекция (корректор)}: полученное прогнозное значение подставляется в неявную формулу (например, метод Адамса-Мултона), и уточнённое значение вычисляется итерационно (обычно 1–2 итерации).
\end{enumerate}
Это позволяет повысить устойчивость и точность, приближаясь к точности неявного метода без решения нелинейных уравнений.

\item \textbf{Когда в методах прогноза и коррекции можно переходить на следующий этап вычислений?}

Переход на следующий шаг допустим, если после итераций коррекции разница между двумя последовательными приближениями стала меньше заданной точности \( \varepsilon \) (или выполнилось одно фиксированное число итераций, например, 2). На практике часто выполняют одну итерацию коррекции, после чего переходят к следующему шагу, что даёт хороший баланс точности и затрат.

\item \textbf{Что такое правило Рунге и как оно используется в данной задаче?}

Правило Рунге служит для оценки погрешности численного метода без знания точного решения. Выполняется два расчёта: с шагом \( h \) (решение \( y_h \)) и с шагом \( h/2 \) (решение \( y_{h/2} \)). Тогда главный член погрешности можно оценить по формуле:
\[
\Delta = \frac{|y_{h/2} - y_h|}{2^p - 1},
\]
где \( p \) — порядок точности метода. В данной лабораторной работе правило Рунге применяется для одношаговых методов (Эйлера, Рунге-Кутты) с целью контроля точности и (при необходимости) автоматического уменьшения шага до достижения требуемой точности \( \varepsilon \).

\item \textbf{Чтобы «запустить» метод Адамса, что необходимо вычислить?}

Метод Адамса является многошаговым, например, для 4-шагового явного метода Адамса-Башфорта нужны значения \( y_0, y_1, y_2, y_3 \). Начальное значение \( y_0 \) задано условием. Остальные \( y_1, y_2, y_3 \) нельзя получить самим методом Адамса. Их вычисляют другим одношаговым методом того же или более высокого порядка точности, например, методом Рунге-Кутты 4-го порядка. Этот этап называется \textbf{разгоном} (стартом) многошагового метода.

\end{enumerate}

\end{document}